% !TeX root = ../../Book4.tex
% 纲要标题：### 14.5 非零高阶微分与隐藏延拓的具体模型
% 纲要位置：计划纲要.md，第 2851 行。

\section{非零高阶微分与隐藏延拓的具体模型}
\label{b4:sec:1405}

% 纲要内容：
%
% * 构造四个非零项 \(C^{0,1}=C^{1,0}=C^{1,1}=C^{2,0}=k\) 的第一象限双复形；取 \(d_v:C^{1,0}\to C^{1,1}\) 及 \(d_h:C^{0,1}\to C^{1,1}\)、\(d_h:C^{1,0}\to C^{2,0}\) 为恒等映射，其余微分为零
% * 按列先取上同调：第1页只剩 \((0,1)\) 与 \((2,0)\) 两项，第1阶微分为零；用代表元提升计算非零的 \(d_2:E_2^{0,1}\to E_2^{2,0}\)
% * 直接核验总复形无上同调，并与第3页消失一致；用这个模型说明“第1阶微分为零”不意味着谱序列已经退化
% * 比较 \(\mathbb Z/4\mathbb Z\) 的滤过 \(0\subset2\mathbb Z/4\mathbb Z\subset\mathbb Z/4\mathbb Z\) 与 \((\mathbb Z/2\mathbb Z)^2\) 的直和因子滤过：伴随分次相同而原群不同
% * 把高阶微分的计算与极限页后的延拓问题明确分开；本节为正文算例，不假定知道代数拓扑中的谱序列
%
% ---
%

\subsection{四项第一象限双复形}
\label{b4:subsec:1405:01}

设 \(k\) 是任意域。只在以下四个位置放置一维向量空间：
\[
 C^{0,1}=ka,\qquad C^{1,0}=kb,\qquad
 C^{1,1}=kc,\qquad C^{2,0}=ke.
\]
定义 \(d_ha=c\)、\(d_vb=c\)、\(d_hb=e\)，其他箭头均为零。这三个非零箭头在所选基下都是恒等映射。每个二次横或纵复合都为零；混合的两步路径也都为零，因此满足 \(d_hd_v+d_vd_h=0\)。这里无需把其中一条箭头改为负号，因为没有一个非零的两步方格。
\[
\begin{tikzcd}[column sep=large,row sep=large]
 ka \arrow[r,"1"] & kc & 0\\
 0 & kb \arrow[u,"1"] \arrow[r,"1"'] & ke
\end{tikzcd}
\]
上排是 \(q=1\)，下排是 \(q=0\)，从左至右横指标为 \(0,1,2\)。这幅图的斜向联系要经过一次求原像才能出现，正适合检验高阶微分。

\subsection{\texorpdfstring{列滤过与非零的 \(d_2\)}{列滤过与非零的 d2}}
\label{b4:subsec:1405:02}

按列滤过总复形，先取纵向上同调。中间列 \(kb\xrightarrow{1}kc\) 正合，故第一页只剩
\[
 E_1^{0,1}=k[a],\qquad E_1^{2,0}=k[e].
\]
没有相邻非零项，\(d_1=0\)，所以 \(E_2=E_1\)。但代表元 \(a\) 不是总余圈：\(da=c\) 尚在滤过一层。为了升到第二页，要加上较深一层的修正 \(-b\)，因为
\[
 d(a-b)=c-(c+e)=-e\in F^2\operatorname{Tot}(C)^2.
\]
于是由页的定义直接得到
\begin{equation}\label{b4:eq:nonzero-d2}
 d_2([a])=-[e].
\end{equation}
该映射在任意域上都是同构，在特征二时负号等于正号。因此 \(E_3=0\)。这个负号来自修正代表元时减去 \(b\)，并非来自再次改变双复形的反交换约定。

\subsection{总复形无上同调性与谱序列退化}
\label{b4:subsec:1405:03}

还可以直接检查总复形。它只有次数一和二，按 \((a,b)\)、\((c,e)\) 排列基，微分矩阵为
\[
 k^2\xrightarrow{\left(\begin{smallmatrix}1&1\\0&1\end{smallmatrix}\right)}k^2.
\]
行列式为一，所以总复形无上同调。这与 \(E_3=E_\infty=0\) 完全吻合。

\ProseParagraphBreak
若在 \(d_1=0\) 后就停止计算，会错误地给总次数一、二各留下一个一维上同调。谱序列的高阶微分正是从“当前分层下看似闭合”到“能否成为真正余圈”的进一步检验。本例中只需一次修正；更长的双复形中会出现多个相继修正。

\subsection{相同伴随分次的不同交换群}
\label{b4:subsec:1405:04}

现在取两个集中在次数零、微分为零的交换群复形。第一个是 \(A=\mathbb Z/4\)，滤过为
\[
 F^0A=A,\qquad F^1A=2A,\qquad F^2A=0.
\]
第二个是 \(B=(\mathbb Z/2)^2\)，取 \(F^0B=B\)、\(F^1B=0\oplus\mathbb Z/2\)、\(F^2B=0\)。两者在 \(p\le0\) 延拓为整个对象，在 \(p\ge2\) 延拓为零。它们的伴随分次均在 \(p=0,1\) 各有一份 \(\mathbb Z/2\)。因为总微分为零，两套谱序列从第零页起完全相同，所有后续微分也都为零。

\ProseParagraphBreak
然而 \(A\) 有四阶元素，\(B\) 的每个元素都被二零化，故二者不同构。对 \(A\)，短正合列 \(0\to2A\to A\to A/2A\to0\) 不分裂：商的非零元若有二阶提升，就能产生分裂，但它的两个提升 \(\overline1,\overline3\) 都是四阶。对 \(B\)，投影到第一因子显然有截面。两种对象的差别正是 \(\operatorname{Ext}_{\mathbb Z}^1(\mathbb Z/2,\mathbb Z/2)\cong\mathbb Z/2\) 中的非零类与零类。

\subsection{高阶微分与极限页后的延拓}
\label{b4:subsec:1405:05}

前一个双复形告诉我们：第一页留下的类可能被第二阶微分消去。后两个滤过群告诉我们：即使全部微分已知且永远为零，仍可能不能识别目标群。前者是代表元能否继续修正的问题，后者是已经得到的相邻商如何拼成原对象的问题。

\ProseParagraphBreak
在实际计算中，先列出各页非零位置，逐一判定可能的箭头；需要高阶微分时回到总复形修正代表；确定稳定页以后，再写出目标的有限滤过及每个相邻商。若最后的扩张没有得到进一步依据，就保留短正合列作为结论。这样既能使用谱序列提供的信息，也不会把它未决定的结构误写成直和。
